%
% CHAPTER - Philosophy
%

\chapterimage{School_of_Athens.pdf}

\chapter{Philosophical Foundations}
\label{chap:Philosophy}

\begin{quote}
\begin{flushright}
\emph{To go where you don't know, \\
you have to go the way you don't know.} \\
San Juan de la Cruz 
\end{flushright}
\end{quote}
\bigskip

The \emph{philosophy of science}\index{Philosophy of science} examines the foundations, methods, and limits of scientific inquiry. Together with related disciplines such as \emph{epistemology}\index{Epistemology}, \emph{ontology}\index{Ontology}, and \emph{metaphysics}\index{Metaphysics}, it addresses fundamental questions concerning the nature of knowledge, the existence and knowability of entities, the interpretation of observations, the construction of scientific representations and explanations, and the justification of scientific hypotheses. These questions provide an essential conceptual foundation for any attempt to characterize scientific knowledge and its limitations mathematically.

The purpose of this chapter is not to provide a general survey of the philosophy of science, but to identify and examine the philosophical problems that a mathematical theory of scientific knowledge and ignorance must confront. Particular attention is given to the fundamental questions, requirements, and alternative approaches relevant to the Theory of Nescience. Where appropriate, we introduce established philosophical problems, discuss competing conceptions, and identify their principal difficulties. The chapter does not attempt to resolve these questions or defend the particular assumptions adopted by the Theory of Nescience; those matters are addressed in the corresponding theoretical chapters. Instead, it establishes the philosophical background against which the theory's definitions, assumptions, and results can be examined.

The chapter is organized around eight closely related subjects. We begin by examining the nature of scientific knowledge and its relationship to belief, information, understanding, and ignorance. We then consider the ontological status of scientific entities and the conditions under which they can be identified and known. Next, we investigate scientific observation and evidence, followed by the problems of scientific representation and the requirements for adequate representations. We subsequently examine the nature of scientific descriptions, models, and explanations, and the criteria used to evaluate model adequacy, accuracy, simplicity, and explanatory power. The chapter then addresses scientific justification, theory choice, and the possibility of characterizing and quantifying scientific progress. Finally, we consider the discovery of new hypotheses, concepts, and entities, including previously unrecognized unknowns, and examine the fundamental limits of scientific knowledge and the possibility of attaining perfect understanding. 

%
% Section: Scientific Knowledge
%

\section{Scientific Knowledge}
\label{sec:scientific_knowledge}

The nature of \emph{scientific knowledge}\index{Scientific knowledge} is a central problem in the philosophy of science and \emph{epistemology}\index{Epistemology}, the branch of philosophy concerned with the nature, acquisition, justification, and limits of knowledge. Although science aims to produce reliable knowledge through observation, reasoning, and systematic investigation, there is no universally accepted definition of knowledge or single method by which it must be acquired. Several fundamental philosophical problems arise when attempting to characterize scientific knowledge and distinguish it from other forms of information, belief, and understanding.

\emph{Problem of Scientific Demarcation}\index{Scientific demarcation}: What distinguishes scientific knowledge from other forms of knowledge, such as mathematical, philosophical, or everyday knowledge? Proposed criteria include empirical testability, falsifiability, reproducibility, and systematic justification. However, no single criterion is universally accepted, and scientific disciplines differ substantially in their methods of inquiry. Furthermore, some legitimate scientific activities, such as theoretical investigations and exploratory research, may not immediately produce testable predictions. The demarcation problem concerns whether necessary and sufficient conditions can be established for identifying scientific knowledge, or whether scientific status depends on a combination of methodological and epistemic characteristics.

\emph{Problem of Knowledge}\index{Problem of knowledge}: What are the necessary and sufficient conditions for something to constitute knowledge? The traditional analysis characterizes propositional knowledge as \emph{justified true belief}. According to this account, knowing a proposition requires believing it, the proposition being true, and possessing adequate justification for that belief. However, the \emph{Gettier problem}\index{Gettier problem} demonstrates that justified true beliefs may arise through epistemic luck without constituting knowledge. Alternative approaches introduce additional conditions involving reliability, causal connections, or intellectual competence. The problem remains whether knowledge admits a universal characterization, particularly when extending the concept beyond individual beliefs to scientific theories, models, and collective knowledge.

\emph{Problem of Objectual Knowledge}\index{Objectual knowledge}: What does it mean to know an entity, rather than merely know that a proposition about it is true? \emph{Propositional knowledge}\index{Propositional knowledge} concerns knowledge that particular statements are true, whereas \emph{objectual knowledge}\index{Objectual knowledge} concerns knowledge of objects, places, phenomena, or other subjects of inquiry. Some philosophical approaches attempt to characterize objectual knowledge in terms of acquaintance or collections of propositions, while others maintain that it constitutes a distinct epistemic relationship. Furthermore, knowing an entity does not necessarily imply knowing all its properties or relationships. The philosophical problem is determining whether knowledge of an entity can be reduced to propositional knowledge, what constitutes a complete account of an entity, and whether the extent of such knowledge can be meaningfully evaluated.

\emph{Knowledge and Information}\index{Knowledge and information}: What distinguishes possessing information from possessing knowledge? Information can be understood syntactically, as symbols or messages, or semantically, as content concerning some subject. In information theory, quantities such as Shannon entropy characterize uncertainty within a specified probabilistic framework without requiring that the information be true, meaningful, or understood. By contrast, semantic accounts of information address meaning and truth, while epistemological accounts of knowledge introduce additional requirements concerning justification or reliability. Consequently, the availability or accumulation of information does not necessarily imply an increase in knowledge. An important question is what conditions allow information to be transformed into knowledge and whether information-theoretic quantities can adequately characterize epistemic properties.

\emph{Knowledge and Understanding}\index{Knowledge and understanding}: Does possessing accurate knowledge necessarily imply understanding? A person may know that a phenomenon occurs, or successfully predict its occurrence, without understanding why it occurs. Scientific understanding is frequently associated with the ability to explain, identify relationships, recognize patterns, or reason about hypothetical situations. Some philosophical accounts regard understanding as a form of knowledge, while others consider it a distinct epistemic achievement involving the grasp of explanatory relationships. This distinction raises the question of whether accurate information, successful prediction, and explanatory understanding should be evaluated according to the same criteria or treated as separate dimensions of scientific knowledge.

\emph{Problem of Epistemic Justification}\index{Epistemic justification}: What makes a scientific belief or hypothesis sufficiently justified to qualify as knowledge? Foundationalist approaches seek to ground justification in basic beliefs or evidence, while coherentist approaches emphasize consistency and mutual support among beliefs. Reliabilist approaches instead focus on the dependability of the processes through which beliefs are formed. These alternatives disagree about the relationship between evidence, reasoning, and knowledge. Moreover, scientific conclusions are generally provisional and subject to revision, raising the question of how knowledge can be justified without requiring absolute certainty.

\emph{Problem of Quantifying Knowledge}\index{Quantification of knowledge}: Can the extent or degree of knowledge be measured? Traditional epistemology frequently treats knowledge as a condition that either holds or does not hold for a particular proposition. However, scientific inquiry involves partial, incomplete, and progressively refined accounts of phenomena. Quantifying knowledge requires identifying which epistemic properties are measurable, establishing suitable standards of comparison, and determining whether different forms of knowledge can be expressed on a common scale. Probabilistic measures, information-theoretic quantities, and measures of explanatory or descriptive complexity provide possible approaches, but none is universally accepted as a general measure of scientific knowledge. In particular, uncertainty, information content, descriptive complexity, and epistemic adequacy are distinct concepts that need not vary together.

\emph{Problem of Defining Ignorance}\index{Ignorance}: What constitutes ignorance, and is ignorance simply the absence of knowledge? The \emph{knowledge view}\index{Knowledge view of ignorance} characterizes ignorance in terms of lacking knowledge, while the \emph{true-belief view}\index{True-belief view of ignorance} associates it with the absence of true belief. Alternative approaches emphasize failures of epistemic access, inquiry, or the capacities through which knowledge is acquired. A further distinction concerns \emph{propositional ignorance}\index{Propositional ignorance}, involving particular facts or questions, and \emph{objectual ignorance}\index{Objectual ignorance}, concerning an entity or subject matter more generally. The latter may admit degrees depending on the extent, organization, and completeness of the knowledge possessed about the subject. These approaches raise the question of whether ignorance is a single epistemic condition or a collection of distinct limitations, and whether its characterization necessarily depends on the adopted conception of knowledge.

\emph{Problem of Quantifying Ignorance}\index{Quantification of ignorance}: Can ignorance be characterized and measured independently of knowledge? Different conceptions of ignorance may lead to different quantitative characterizations, depending on whether the relevant limitations concern knowledge, belief, epistemic access, or understanding. A distinction must be made between ignorance concerning an identified subject and ignorance concerning subjects whose existence or significance has not yet been recognized. Furthermore, any quantitative measure must specify the domain of inquiry, the standards against which ignorance is evaluated, and the conditions under which different states of ignorance can be compared. It remains an open philosophical question whether ignorance admits a general quantitative characterization, or whether distinct forms of ignorance require different measures.

%
% Section: Ontology and Scientific Entities
%

\section{Ontology and Scientific Entities}
\label{sec:what-is-an-entity}

\emph{Ontology}\index{Ontology} is the branch of philosophy concerned with existence and the nature of entities. Scientific inquiry investigates a wide variety of entities, including physical objects, phenomena, processes, mathematical structures, and theoretical constructs. However, determining what constitutes an entity, whether an entity exists, and under what conditions it can be known raises several fundamental philosophical problems. These questions are particularly important when establishing the domain of scientific investigation and the conditions under which entities can be represented, distinguished, classified, and understood.

\emph{Problem of Existence}\index{Problem of existence}: What does it mean for an entity to exist, and what kinds of entities should be admitted into scientific inquiry? Ontological realism generally maintains that reality exists independently of our knowledge or conceptualization of it, whereas idealist and constructivist approaches assign different roles to minds, concepts, or practices in constituting the objects of investigation. A related issue is that of \emph{ontological commitment}\index{Ontological commitment}: under what conditions does a scientific theory commit us to the existence of the entities it postulates? The use of an entity in a scientific explanation does not necessarily establish its existence. This raises the question of whether scientific investigation should be restricted to entities whose existence has been established or should also accommodate entities whose existence remains hypothetical.

\emph{Problem of Identity and Individuation}\index{Identity and individuation}: What makes an entity distinct from another, and under what conditions can two descriptions be said to refer to the same entity? The problem of identity concerns the conditions under which an entity remains the same despite changes in its properties, state, or representation. The related problem of individuation concerns how entities are distinguished and counted. Different criteria may be appropriate for physical objects, processes, mathematical structures, or social institutions. For example, two physical objects may share all their measured properties without being numerically identical, while two mathematical expressions may denote the same object despite having different forms. Establishing identity criteria is essential for determining when apparently different scientific accounts concern the same entity.

\emph{Problem of Abstract Entities}\index{Abstract entities}: Do abstract entities, such as numbers, mathematical structures, and algorithms, exist independently of human thought? \emph{Mathematical Platonism}\index{Mathematical Platonism} maintains that mathematical objects exist independently of our knowledge of them, whereas nominalist approaches deny that abstract objects must be admitted into our ontology. Other approaches, such as structuralism, emphasize mathematical structures and relationships rather than independently existing individual objects. The problem is particularly relevant because abstract entities cannot be investigated through ordinary physical observation, yet they are central to mathematics and theoretical science. An important question is whether the same ontological and epistemological criteria can apply to concrete and abstract entities.

\emph{Problem of Hypothetical and Fictitious Entities}\index{Hypothetical entities}\index{Fictitious entities}: Can scientific inquiry meaningfully investigate entities whose existence has not been established, or which do not exist? Scientific theories frequently introduce hypothetical entities to explain observations, and some of these entities are subsequently rejected. The historical concepts of phlogiston and the luminiferous ether illustrate how scientific investigation can proceed using entities that are later abandoned. This raises the philosophical \emph{problem of reference}\index{Problem of reference}: how can statements about nonexistent entities be meaningful, and what do their descriptions refer to? It also raises the question of whether a hypothetical entity should be distinguished from a fictitious one, and whether the latter can legitimately serve as an object of scientific investigation.

\emph{Scientific Realism and Anti-Realism}\index{Scientific realism}\index{Scientific anti-realism}: Do successful scientific theories describe entities that actually exist, or are they merely useful instruments for organizing and predicting observations? Scientific realism generally maintains that our best scientific theories provide approximately true accounts of a mind-independent reality, including its unobservable entities. Anti-realist approaches, such as constructive empiricism, maintain that accepting a scientific theory need not involve believing in the existence of all the entities it postulates, but may instead require accepting its empirical adequacy. The disagreement is particularly significant for theoretical entities that cannot be independently observed. The central question is whether explanatory and predictive success provides sufficient grounds for believing in the existence of the entities postulated by a theory.

\emph{Problem of Observability}\index{Problem of observability}: What distinguishes observable from unobservable entities, and does this distinction establish a meaningful boundary for scientific knowledge? Some entities can be perceived directly, while others can only be investigated through instruments, indirect measurements, or theoretical inference. However, the distinction between direct observation and indirect detection is not always clear, since scientific instruments and observation procedures frequently depend on theoretical assumptions. Furthermore, observability may change as technology advances. This raises the question of whether observability is an intrinsic property of an entity or a relation between an entity and the available methods of investigation. More fundamentally, it questions whether observability is necessary for knowledge or merely one possible means of acquiring it.

\emph{Problem of Classification and Natural Kinds}\index{Natural kinds}: Do scientific classifications reflect objective divisions in reality, or are they conceptual systems introduced to organize knowledge? The theory of \emph{natural kinds}\index{Natural kinds} proposes that certain categories correspond to genuine groupings in nature, while conventionalist approaches emphasize the role of human purposes and classification practices. Essentialist accounts associate natural kinds with defining properties, whereas alternative approaches allow categories to be characterized by clusters of properties, causal mechanisms, or historical relationships. Since the same entities may be classified in multiple ways, an important problem is determining which classifications are scientifically meaningful. Related questions concern whether there exists a uniquely correct classification and whether the boundaries between scientific categories are necessarily precise.

\emph{Problem of Composition and Emergence}\index{Composition}\index{Emergence}: Should complex systems be considered entities in their own right, or can they be understood entirely through their constituent parts? \emph{Reductionism}\index{Reductionism} emphasizes the possibility of explaining complex phenomena in terms of more fundamental components and their interactions. \emph{Holism}\index{Holism} emphasizes the importance of properties and relationships that arise at the level of the system as a whole. Related debates concern whether emergent properties are reducible to lower-level descriptions. These perspectives raise questions about the appropriate level at which scientific entities should be identified and whether different levels of organization require distinct ontological categories.

\emph{Problem of Knowability}\index{Problem of knowability}: Under what conditions can an entity be known, and are there entities that are fundamentally unknowable? A distinction must be made between entities that are currently unknown, entities that cannot be investigated using available methods, and entities that cannot be known even in principle. These distinctions depend on how knowledge and the permissible means of acquiring it are defined. An entity may be unknowable relative to a particular observational, conceptual, or computational framework while remaining knowable under another. More fundamental limitations may arise from logical, mathematical, or physical constraints. However, establishing that an entity is unknowable in principle requires stronger arguments than demonstrating the limitations of existing methods. Furthermore, unknowability must be distinguished from nonexistence: an entity may exist without being knowable, while a nonexistent hypothetical entity may still be the subject of meaningful scientific investigation. The central problem is whether general criteria can be established for determining the boundaries of possible knowledge.

%
% Scientific Observation and Evidence
%

\section{Scientific Observation and Evidence}
\label{sec:scientific_observation}

\emph{Scientific observation}\index{Scientific observation} provides an essential source of empirical evidence for investigating the natural world. Observations may involve direct sensory perception, scientific instruments, or experimental procedures designed to detect and measure particular phenomena. However, observations do not necessarily provide direct or complete access to reality. Their production and interpretation depend on measurement procedures, conceptual assumptions, and methodological choices. Several fundamental philosophical problems arise when determining how observations become scientific evidence and under what conditions that evidence can be considered reliable, objective, and relevant.

\emph{Problem of Observational Access}\index{Observational access}: Under what conditions can information about an entity or phenomenon be obtained through observation? Traditional empiricist approaches emphasize sensory experience as a foundation of knowledge, while contemporary scientific practice relies extensively on instruments that transform physical interactions into observable records. Such instruments may provide indirect access to phenomena that cannot be perceived directly. However, interpreting their outputs requires assumptions about the relationship between the instrument and the phenomenon under investigation. The central problem is determining when an observable signal provides legitimate information about its intended target, particularly when the target cannot be independently observed.

\emph{Theory-Ladenness of Observation}\index{Theory-ladenness}: Can observations provide evidence independently of the theories they are intended to evaluate? The thesis of theory-ladenness maintains that observations are influenced by prior knowledge, conceptual frameworks, and theoretical expectations. These influences may affect what researchers perceive, which observations they consider relevant, and how they interpret experimental results. Strong versions of this thesis challenge the possibility of theory-neutral evidence, while more moderate accounts recognize that theoretical assumptions do not necessarily invalidate observations. The principal difficulty concerns the potential circularity of using observations that depend on a theory to justify that same theory. An important question is whether independent procedures, alternative interpretations, or comparisons across different theoretical frameworks can overcome this difficulty.

\emph{Problem of Measurement}\index{Scientific measurement}: How can properties of scientific entities be translated into observable and measurable quantities? Measurement requires establishing a relationship between a property of interest and the values or categories assigned through a measurement procedure. Operationalist approaches emphasize the procedures through which quantities are measured, whereas representational and model-based approaches examine how measurement results correspond to properties and relationships in the world. Different procedures may yield different results, and the validity of a measurement may depend on assumptions about the instruments, the target, and the relevant physical or conceptual framework. This raises the problem of establishing measurement validity and determining whether distinct procedures measure the same property. A related difficulty concerns calibration: how can the accuracy of a measurement procedure be established when the standards used to evaluate it are themselves fallible?

\emph{Data and Phenomena}\index{Data and phenomena}: What distinguishes observational data from the phenomena that scientific theories seek to explain? The \emph{data--phenomena distinction}, developed by Bogen and Woodward, emphasizes that individual experimental records are not necessarily identical to the stable phenomena inferred from them. Observational data may contain noise, measurement artifacts, and variations associated with particular experimental conditions, while scientific theories often seek to explain more general patterns or regularities. Establishing these patterns may require statistical analysis, data processing, and additional assumptions. The philosophical problem concerns how reliable claims about phenomena can be inferred from imperfect observational records, and whether the distinction between data and phenomena can always be maintained.

\emph{Problem of Observational Reliability}\index{Observational reliability}: What makes an observation or experimental result sufficiently reliable to serve as scientific evidence? Observations may be affected by random errors, systematic biases, instrument limitations, environmental disturbances, or incorrect interpretations. Measurement precision must therefore be distinguished from accuracy: repeated measurements may produce consistent results while remaining systematically incorrect. Approaches to establishing reliability include calibration, error analysis, repeated measurements, independent instruments, and methodological controls. However, these procedures may themselves depend on assumptions that cannot be established with absolute certainty. The central problem is determining the conditions under which observational errors can be identified, quantified, or sufficiently controlled, particularly when an independent reference is unavailable.

\emph{Problem of Scientific Objectivity}\index{Scientific objectivity}: Can observational evidence be objective when its production and interpretation depend on human decisions and theoretical assumptions? Traditional conceptions of objectivity emphasize the independence of observations from individual beliefs and expectations. Alternative approaches emphasize intersubjective agreement, reproducibility, methodological transparency, and independent critical scrutiny. These approaches recognize that objectivity need not require the complete absence of theoretical assumptions, but may instead depend on procedures for identifying and correcting their influence. Nevertheless, reproducibility alone does not guarantee correctness, since independent observers may share the same methodological limitations or systematic errors. The problem is to determine which conditions allow observational evidence to support claims independently of the particular individuals or procedures involved in its production.

\emph{Problem of Evidential Relevance}\index{Evidential relevance}: What makes an observation relevant to a particular scientific question or hypothesis? Scientific inquiry requires selecting which entities, properties, and phenomena to investigate, since observing every potentially relevant aspect of reality is generally impossible. The relevance of evidence depends partly on the question under investigation and on assumptions concerning which observations could support or challenge a proposed explanation. However, selecting observations according to existing expectations may introduce selection bias or conceal unexpected phenomena. Controlled experiments attempt to isolate relevant factors through intervention, but their design also depends on assumptions about which variables should be manipulated or controlled. The philosophical problem is determining how relevant evidence can be identified without excluding information that may challenge established knowledge.

\emph{Problem of Observational Incompleteness}\index{Observational incompleteness}: Can a finite collection of observations provide sufficient information about an entity or phenomenon? Scientific observations are necessarily constrained by the properties measured, the resolution of available instruments, the conditions under which measurements are obtained, and the extent of the observed sample. Consequently, different entities or processes may produce indistinguishable observational records under a given procedure, even when they differ in important respects. Furthermore, the absence of an observed phenomenon does not necessarily establish its absence, particularly when detection capabilities are limited. These difficulties raise questions about whether observational completeness can be meaningfully defined, how missing information can be identified, and under what conditions additional observations can distinguish between otherwise observationally equivalent possibilities.

%
% Scientific Representation
%

\section{Scientific Representation}
\label{sec:scientific_representation}

\emph{Scientific representation}\index{Scientific representation} concerns the relationship between the objects used in scientific inquiry and the entities or phenomena they are intended to represent. Scientific representations take many forms, including measurements, diagrams, images, datasets, mathematical structures, and theoretical models. They allow researchers to investigate entities that may not be directly accessible, communicate information about them, and reason about their properties and behavior. However, the nature of the representational relationship remains philosophically contested. In particular, it is necessary to distinguish between the conditions under which something represents an entity and the conditions under which it represents that entity accurately or adequately. Although much of the philosophical literature focuses on scientific models, the problems discussed below concern scientific representation more generally.

The principal philosophical problems concerning scientific representation are the following:

\emph{Scientific Representation Problem}\index{Scientific representation problem}: What are the necessary and sufficient conditions for an object to represent a scientific entity? The problem concerns the relationship between a \emph{representational vehicle}, such as a diagram or model, and its \emph{target}, the entity or phenomenon being represented. A central difficulty is determining what establishes this relationship: similarity, structural correspondence, interpretation, inferential use, or some other condition. Moreover, being a representation does not necessarily imply being an accurate representation. A scientific model may represent its intended target while providing an incomplete or incorrect account of it. The problem is therefore to explain what establishes a representational relationship independently of the degree to which the representation succeeds.

\emph{Representational Demarcation Problem}\index{Representational demarcation}: Are scientific representations fundamentally different from other forms of representation, such as those used in art, language, or everyday communication? Some philosophical approaches seek distinctive conditions for scientific representation, emphasizing its role in producing knowledge about a target. Others maintain that scientific representations do not constitute a separate category and that the same general principles of representation apply across different domains. The difficulty is determining whether the epistemic requirements associated with scientific practice establish a distinctive kind of representation or merely impose additional standards on representations in general.

\emph{Problem of Style}\index{Problem of style}: What different styles of scientific representation are possible, and how can they be characterized? The same entity may be represented through measurements, diagrams, mathematical structures, or other representational systems, each emphasizing different aspects of its target. Representational styles differ not only in their physical or symbolic form but also in the conventions, interpretations, and relationships they establish between representations and their targets. Furthermore, new styles may emerge as scientific methods evolve. The problem is to determine whether representational styles can be systematically classified and how differences between them affect the information conveyed about an entity.

\emph{Standards of Accuracy}\index{Standards of accuracy}: What makes a scientific representation accurate, and how should its accuracy be evaluated? Representations typically capture selected aspects of their targets while omitting, simplifying, or idealizing others. Consequently, representational accuracy cannot always be understood as complete correspondence with reality. Different research purposes may require different standards, and a representation that is adequate for one purpose may be inadequate for another. The problem is to establish criteria for comparing representations, distinguishing acceptable approximations from errors, and determining whether accuracy can be evaluated independently of the intended use of the representation.

\emph{Problem of Ontology}\index{Ontology of representations}: What kinds of objects can serve as scientific representations? Representational vehicles may be concrete objects, such as physical models and scientific instruments, or abstract objects, such as mathematical structures and theoretical constructs. This raises questions about whether representations possess their representational properties intrinsically or acquire them through interpretation and scientific practice. In particular, abstract mathematical objects do not ordinarily refer to physical entities merely by existing. The problem is to determine what ontological conditions, if any, must be satisfied by a representational vehicle and how objects of different kinds can fulfill the same epistemic function.

Beyond these general problems, philosophical accounts of scientific representation are commonly evaluated against several \emph{conditions of adequacy}\index{Conditions of adequacy}. These conditions concern the capacities and characteristics that an account of representation should explain, rather than properties that every individual representation must necessarily possess.

\emph{Requirement of Directionality}\index{Requirement of directionality}: Scientific representation is generally asymmetric: a representation is about its target, but the target does not thereby represent the representation. For example, a map represents a geographical region, whereas the region does not ordinarily represent the map. This asymmetry cannot be explained solely through symmetric relationships such as similarity or isomorphism. An adequate account must therefore explain what establishes the direction of the representational relationship, whether through reference, interpretation, intentionality, or other mechanisms.

\emph{Surrogative Reasoning}\index{Surrogative reasoning}: Scientific representations allow researchers to investigate a target indirectly by reasoning about its representation. A mathematical model, for example, may support inferences about the behavior of a physical system without requiring direct experimentation on that system. However, not every inference concerning a representation is necessarily valid for its target. An adequate philosophical account must explain how conclusions obtained from representations can legitimately be transferred to the entities they represent, and under what conditions such inferences are reliable.

\emph{Applicability of Mathematics}\index{Applicability of mathematics}: How can abstract mathematical structures represent concrete entities and physical phenomena? Mathematics provides structures, relationships, and operations that can be used to investigate scientific systems, yet mathematical objects do not inherently possess the physical properties of their targets. The problem is to explain how mathematical relationships acquire empirical significance and why reasoning within an abstract mathematical framework can produce valid conclusions about the world. This requires accounting for the interpretation and correspondence between mathematical structures and the properties of the entities they represent.

\emph{Possibility of Misrepresentation}\index{Possibility of misrepresentation}: Can something represent an entity even when it represents that entity incorrectly? Scientific representations may omit relevant properties, introduce nonexistent features, or otherwise provide inaccurate accounts of their targets. Nevertheless, such representations can remain meaningful and scientifically useful. An adequate philosophical account must therefore distinguish between failing to represent an entity and representing it inaccurately. This distinction is particularly important because any conception that requires perfect correspondence between representation and target would exclude the possibility of misrepresentation by definition.

\emph{Targetless Models}\index{Targetless models}: Can a scientific representation exist without an actual target? Scientists frequently develop models of hypothetical systems, idealized objects, or entities whose existence has not been established. Some models describe physically impossible situations, while others concern systems that might exist but have never been observed. These cases challenge accounts that require an existing target as a necessary condition for representation. An adequate philosophical account must explain whether such models qualify as representations, what their representational content consists of, and how they can support scientific reasoning despite lacking an actual referent.

Several philosophical approaches have been proposed to explain scientific representation and address these problems. The following accounts illustrate the principal alternatives, together with their main conceptual difficulties.

\emph{Stipulative Fiat}\index{Stipulative fiat}: According to the stipulative approach, a scientific model $M$ represents a target $T$ if a user stipulates that $M$ represents $T$. Representation is therefore established through intentional designation rather than an intrinsic relationship between the model and its target. This approach accommodates considerable freedom in the choice of representational vehicles. However, stipulation alone does not explain how the representation supports reliable inferences about its target. The principal difficulty is distinguishing merely assigning representational status from establishing the epistemic usefulness of a representation.

\emph{Similarity Conception}\index{Similarity conception}: Similarity-based approaches maintain that a scientific model $M$ represents a target $T$ in virtue of relevant similarities between them. Shared properties may provide a basis for transferring conclusions from the model to the target, thereby supporting surrogative reasoning. However, similarity is generally symmetric, whereas representation is directional. Furthermore, two objects may resemble one another without either representing the other. The approach must therefore specify which similarities are relevant, how they are selected, and why those similarities establish a representational relationship rather than mere resemblance.

\emph{Structuralist Conception}\index{Structuralist conception}: Structuralist approaches explain scientific representation through mathematical relationships between structures associated with models and their targets. A particularly strong formulation requires an isomorphism between the relevant structures, while more flexible approaches consider partial isomorphisms, homomorphisms, or other structural correspondences. Such relationships can support mathematical reasoning about the target, but structural correspondence alone does not establish which entity is being represented or why the correspondence is scientifically meaningful. Moreover, exact isomorphism may be too restrictive for approximate or idealized representations. The main difficulties concern the interpretation, directionality, and adequacy of the structural relationship.

\emph{Inferential Conception}\index{Inferential conception}: Inferential approaches emphasize the capacity of a representation to support reasoning about its target. In particular, Suárez's inferential conception identifies representational directionality and the capacity for surrogative inference as fundamental conditions of scientific representation. Rather than requiring a universal relationship of similarity or structural correspondence, this approach emphasizes the role of representational practices and the inferences they make possible. Its principal challenge is explaining what makes such inferences legitimate and how representational content and accuracy can be established without relying on the very relationships the approach seeks to avoid treating as fundamental.

\emph{Fiction View of Models}\index{Fiction view of models}: Fiction-based approaches understand scientific models through their relationship with imagination and make-believe. Models may prescribe imagining certain properties or situations, allowing researchers to investigate possibilities without requiring the modeled systems to exist exactly as described. This perspective accommodates idealization, abstraction, and hypothetical entities. However, it raises questions about how imagined properties relate to actual targets and how reasoning based on fictional constructions can produce reliable scientific knowledge. A further difficulty concerns the ontological status of the imagined systems themselves.

\emph{Representation-As}\index{Representation-as}: The representation-as approach emphasizes that scientific representations depict their targets as possessing particular properties or structures. Rather than reproducing every aspect of an entity, a representation attributes selected characteristics to it through interpretation, denotation, and the association of representational properties with target properties. This approach accommodates differences in representational style and allows representations to be inaccurate without losing their representational status. However, it must explain how the relevant properties are selected, how their interpretation is established, and under what conditions attributing them to a target supports justified scientific inferences.

These approaches illustrate the diversity of possible accounts of scientific representation. No single conception commands universal acceptance, and different approaches emphasize different aspects of the relationship between representations and their targets. The central philosophical challenge remains to explain how representations acquire meaning, how their adequacy can be evaluated, and how they permit reliable reasoning about entities that may be only partially accessible through scientific investigation.

%
% Scientific Description and Explanation
%

\section{Scientific Description and Explanation}
\label{sec:scientific_explanation}

\emph{Scientific description}\index{Scientific description} and \emph{scientific explanation}\index{Scientific explanation} are fundamental activities of scientific inquiry. Scientists develop descriptions, models, theories, and laws to characterize phenomena, identify regularities, generate predictions, and improve their understanding of the world. However, these concepts are not equivalent, and their relationships remain subjects of philosophical debate. In particular, an account may accurately describe or predict a phenomenon without explaining why it occurs. Several fundamental problems arise when attempting to characterize scientific explanations and establish the conditions under which they contribute to scientific understanding.

\emph{Problem of Scientific Description}\index{Scientific description}: What constitutes a scientific description, and how does it differ from a model, theory, or law? A scientific description characterizes properties, relationships, or behavior associated with an entity or phenomenon. A \emph{scientific model}\index{Scientific model} typically provides a simplified or idealized account of some aspect of a target, often expressed through mathematical structures, diagrams, or computational procedures. A \emph{scientific theory}\index{Scientific theory} generally provides a broader conceptual or mathematical framework that organizes and relates scientific claims, while a \emph{scientific law}\index{Scientific law} expresses a regularity or relationship regarded as holding under specified conditions. However, these distinctions are not universally accepted, and models may function independently of comprehensive theories. The philosophical problem is determining whether these categories admit precise and general definitions, how they relate to one another, and which characteristics distinguish their epistemic functions.

\emph{Problem of Scientific Explanation}\index{Problem of scientific explanation}: What distinguishes an explanation from a description of a phenomenon? Scientific explanations are commonly understood as accounts that answer questions about why or how something occurs, rather than merely stating what occurs. However, a description may be entirely accurate without providing an explanation. For example, identifying a regularity in observational data does not necessarily establish why that regularity exists. Proposed criteria for explanation include deduction from laws, identification of causal mechanisms, statistical relevance, and the unification of phenomena under general principles. The central problem is determining what additional relationship, beyond descriptive correctness, makes an account explanatory.

\emph{Problem of Explanatory Adequacy}\index{Explanatory adequacy}: What conditions must an explanation satisfy to be considered scientifically adequate? Proposed requirements include empirical adequacy, logical consistency, explanatory relevance, generality, and the capacity to support appropriate inferences. However, these characteristics may not be jointly sufficient for explanation. A logically consistent account may fit the observations while failing to identify the factors responsible for the phenomenon. Furthermore, an explanation may be adequate for one scientific question while remaining incomplete for another. The problem is determining whether explanatory adequacy can be established through universal criteria or whether it necessarily depends on the phenomenon, explanatory purpose, and available background knowledge.

\emph{Explanation and Prediction}\index{Explanation and prediction}: Is explaining a phenomenon equivalent to predicting it? The \emph{symmetry thesis}\index{Symmetry thesis}, associated with the deductive-nomological account, maintains that explanations and predictions share the same logical structure, differing primarily in whether the phenomenon is already known. However, successful predictions do not necessarily constitute explanations. For example, the height of a flagpole can be used to predict the length of its shadow, and the shadow length can conversely be used to calculate the flagpole's height under suitable conditions. Yet these relationships do not ordinarily provide equally satisfactory explanations in both directions. This difficulty raises the question of what distinguishes explanatory relevance from predictive success and whether explanation requires additional causal, structural, or contextual relationships.

\emph{Problem of Causal Explanation}\index{Causal explanation}: Must a scientific explanation identify the causes responsible for a phenomenon? Causal approaches maintain that explaining an event involves identifying the processes, mechanisms, or dependencies responsible for its occurrence. However, statistical relationships do not necessarily establish causation, and different causal structures may be compatible with the same observational evidence. Moreover, some scientific explanations appeal to mathematical constraints, symmetries, or structural properties rather than causal processes. A related issue is \emph{explanatory asymmetry}\index{Explanatory asymmetry}: the fact that the relationship between what explains and what is explained is often directional. The philosophical problem is determining whether causality is necessary for scientific explanation, how causal relevance can be established, and whether a general account can accommodate both causal and non-causal explanations.

\emph{Problem of Laws and Generality}\index{Scientific laws}\index{Generality}: What role do scientific laws and general principles play in explanation? Some philosophical accounts maintain that explanations require subsuming particular phenomena under general laws. However, distinguishing genuine laws from accidental regularities is itself a longstanding philosophical problem. Furthermore, scientific generalizations may depend on idealized assumptions or hold only within restricted domains. Explanations may also differ in scope, accounting for individual events, classes of phenomena, or relationships across multiple disciplines. The problem is determining whether generality is necessary for explanation, how explanatory scope should be evaluated, and whether explanations based on general principles are inherently preferable to those addressing particular cases.

\emph{Problem of Explanatory Unification}\index{Explanatory unification}: Does scientific understanding improve when multiple phenomena can be explained through a common set of principles? Unificationist approaches emphasize the reduction of independent explanatory patterns and the integration of apparently unrelated phenomena within a common framework. For example, Newtonian mechanics provides a unified account of terrestrial and celestial motion. However, unification does not necessarily establish that the resulting explanation captures the causal or structural relationships responsible for the phenomena. Furthermore, different phenomena may require different explanatory frameworks, and greater generality may come at the expense of important distinctions. The philosophical problem is determining whether unification is a fundamental requirement of scientific explanation or one explanatory virtue among several.

\emph{Explanation and Scientific Understanding}\index{Scientific understanding}: What is the relationship between possessing an explanation and understanding the phenomenon it explains? Scientific understanding is frequently associated with grasping relationships, recognizing regularities, identifying dependencies, and reasoning about alternative possibilities. However, an explanation may be formally correct without being comprehensible to a particular researcher, while a simplified or idealized account may promote understanding without describing every relevant detail accurately. Philosophical approaches differ on whether understanding is primarily an objective property of explanatory relationships or also depends on the cognitive capacities and background knowledge of the individual. The problem is determining whether understanding can be characterized independently of explanation and whether the explanatory value of a model can be evaluated without reference to the person using it.

Several philosophical approaches have been proposed to characterize scientific explanation and address these problems. The following accounts illustrate the principal alternatives and their main limitations.

\emph{Deductive-Nomological Model}\index{Deductive-nomological model}: The deductive-nomological (DN) model, associated with Hempel and Oppenheim, explains a phenomenon by deriving it logically from general laws and relevant conditions. An explanation consists of the \emph{explanans}\index{Explanans}, containing the explanatory premises, and the \emph{explanandum}\index{Explanandum}, describing the phenomenon to be explained. The explanandum must follow deductively from the explanans, which must include at least one law essential to the derivation. This approach provides a precise logical account of explanation and establishes a close relationship between explanation and prediction. However, logical derivability alone does not guarantee explanatory relevance, and the model has difficulty accounting for explanatory asymmetry, probabilistic phenomena, and explanations that do not depend on general laws.

\emph{Statistical-Relevance Model}\index{Statistical-relevance model}: The statistical-relevance (SR) model, developed by Salmon, emphasizes statistically relevant relationships rather than strict logical derivation from deterministic laws. Within a reference class $A$, a factor $C$ is statistically relevant to an outcome $B$ when
\[
P(B \mid A,C) \neq P(B \mid A),
\]
provided the conditional probabilities are defined. This approach accommodates probabilistic phenomena and identifies factors that change the probability of an outcome. However, statistical relevance alone does not establish explanatory or causal relevance. Correlations may arise from common causes, selection effects, or other relationships that do not directly explain the phenomenon. The principal difficulty is determining which statistical relationships are genuinely explanatory and how an appropriate reference class should be selected.

\emph{Causal-Mechanical Accounts}\index{Causal-mechanical explanation}: Causal-mechanical approaches explain phenomena by identifying the causal processes and mechanisms responsible for their occurrence. Rather than relying exclusively on logical derivation or statistical association, these accounts emphasize the relationships through which causes produce effects. Related interventionist approaches, such as Woodward's account, characterize causal explanation through dependencies that remain informative under suitable interventions or counterfactual changes. These approaches address explanatory relevance and asymmetry more directly than purely logical accounts. However, identifying causal mechanisms may be difficult, and not every scientifically valuable explanation appears to require a causal interpretation. Their principal challenge is determining the scope and conditions under which causal relationships provide an adequate basis for explanation.

\emph{Unificationist Accounts}\index{Unificationist explanation}: Unificationist approaches, notably developed by Friedman and Kitcher, understand explanation as the integration of multiple phenomena under a comparatively small collection of general principles or argument patterns. An explanation contributes to understanding by reducing the number of independent facts or assumptions required to account for the phenomena under investigation. This approach emphasizes the relationship between generality, explanatory scope, and scientific understanding. However, a highly unified account need not identify the causal mechanisms responsible for particular phenomena, and different ways of organizing the same body of knowledge may yield different degrees of unification. The central difficulty is determining which forms of unification are genuinely explanatory rather than merely economical or systematic.

\emph{Pragmatic Accounts}\index{Pragmatic explanation}: Pragmatic approaches emphasize the context-dependent character of scientific explanations. According to these accounts, an explanation is an answer to a question posed within a particular context, and its adequacy depends partly on the explanatory interests and background assumptions of the inquirer. For example, explaining why a bridge collapsed may require different information depending on whether the question concerns structural failure, environmental conditions, or maintenance decisions. This perspective accommodates the diversity of explanatory practices across scientific disciplines and research purposes. However, it raises the question of whether explanatory adequacy can be evaluated independently of the interests and interpretations of particular individuals. The principal difficulty is reconciling contextual explanatory requirements with objective standards of scientific knowledge.

These approaches emphasize different characteristics of scientific explanation, including logical derivation, statistical relevance, causality, unification, and contextual adequacy. No single account is universally accepted, and the different approaches need not be mutually exclusive. The central philosophical challenge remains to determine what distinguishes explanation from description and prediction, what makes explanations scientifically adequate, and under what conditions they contribute to genuine understanding.

%
% Section: Model Adequacy and Simplicity
%

\section{Model Adequacy and Simplicity}
\label{sec:model_adequacy}

Scientific models are evaluated according to several criteria, including their ability to account for observations, generate reliable predictions, provide explanations, and capture relevant features of the phenomena under investigation. However, models that account for the same evidence may differ substantially in their assumptions, complexity, and explanatory power. A highly complex model may reproduce observations accurately without revealing meaningful regularities, while a simpler model may provide useful insights despite introducing approximations or idealizations. Determining what constitutes an adequate scientific model therefore requires examining the relationships between accuracy, simplicity, complexity, and explanatory value.

\emph{Problem of Model Adequacy}\index{Model adequacy}: What makes a scientific model adequate for its intended purpose? Possible criteria include empirical accuracy, predictive reliability, explanatory power, logical consistency, and agreement with established knowledge. However, these criteria are not necessarily equivalent, and a model may satisfy some while failing others. A model designed to predict a phenomenon may be evaluated differently from one intended to explain its underlying mechanisms. Furthermore, adequacy may depend on the properties being investigated, the precision required, and the domain in which the model is applied. The fundamental problem is determining whether universal standards of model adequacy can be established or whether adequacy must always be evaluated relative to particular scientific objectives.

\emph{Problem of Empirical Accuracy}\index{Empirical accuracy}: How accurately must a model account for observations to be considered scientifically adequate? Empirical accuracy concerns the correspondence between a model's predictions or outputs and the available observational evidence. However, observations may contain measurement errors, noise, or incomplete information, making exact agreement neither always possible nor necessarily desirable. Different measures of error may also lead to different evaluations of the same model. Furthermore, a model may reproduce existing observations accurately while failing to account for new observations. The problem is determining appropriate standards of accuracy, distinguishing meaningful discrepancies from observational uncertainty, and establishing whether empirical accuracy can be evaluated independently of the intended purpose of the model.

\emph{Problem of Simplicity and Parsimony}\index{Simplicity}\index{Parsimony}: Why should simpler scientific models be preferred over more complex alternatives? The principle of \emph{parsimony}, commonly associated with Occam's razor, recommends avoiding unnecessary assumptions or complexity. However, simplicity admits several interpretations, including fewer parameters, fewer assumptions, shorter descriptions, simpler mathematical expressions, or fewer independent explanatory principles. These criteria need not agree, and a model may be simple according to one criterion while remaining complex according to another. Furthermore, simplicity does not guarantee truth, accuracy, or explanatory adequacy. The philosophical problem is determining whether simplicity has intrinsic epistemic value, whether it can be measured objectively, and under what conditions simpler models should be preferred.

\emph{Problem of Model Complexity}\index{Model complexity}: What constitutes the complexity of a scientific model, and how can it be measured? Possible approaches include counting parameters, measuring the computational resources required to evaluate the model, or determining the length of its shortest description. Information-theoretic approaches, particularly those based on algorithmic complexity, associate complexity with the amount of information required to specify a model. However, complexity may depend on the language, encoding conventions, or computational framework used to describe it. Moreover, a model may contain redundant assumptions, unnecessary parameters, or information that contributes nothing to its explanatory or predictive performance. The problem is distinguishing essential complexity from unnecessary complexity and establishing whether model complexity can be characterized independently of its particular representation.

\emph{Problem of Overfitting and Generalization}\index{Overfitting}\index{Generalization}: Does accurately accounting for existing observations imply that a model will perform well on observations not previously considered? \emph{Overfitting}\index{Overfitting} occurs when a model captures peculiarities of the available data that do not generalize adequately to the underlying phenomenon or to new observations. Conversely, \emph{underfitting}\index{Underfitting} occurs when a model is insufficiently expressive to capture relevant regularities. Statistical learning theory investigates these problems through concepts such as generalization error, model capacity, and the bias--variance trade-off. However, model complexity alone does not determine generalization performance, which also depends on the data, learning procedure, and underlying assumptions. The philosophical problem is determining when agreement with existing evidence provides grounds for expecting reliable performance beyond that evidence.

\emph{Problem of Idealization and Abstraction}\index{Idealization}\index{Abstraction}: Can a scientific model be adequate even when it deliberately omits or misrepresents aspects of reality? \emph{Abstraction} involves disregarding certain features of a target, while \emph{idealization} often involves introducing simplifying assumptions that are not strictly true, such as frictionless surfaces or perfectly rational agents. These practices may produce models that are easier to analyze and that reveal important relationships despite departing from reality. However, simplifications can also conceal relevant mechanisms or lead to misleading conclusions. The problem is determining which omissions and idealizations are legitimate, how their effects on model adequacy should be assessed, and whether greater realism necessarily produces better scientific models.

\emph{Problem of Accuracy--Complexity Trade-offs}\index{Accuracy--complexity trade-off}: How should empirical accuracy be balanced against model complexity? Increasing model complexity may improve agreement with existing observations while introducing additional assumptions, parameters, or descriptive information. Conversely, simplifying a model may reduce complexity at the expense of accuracy. Several approaches address this problem by combining measures of empirical fit with penalties for complexity. Statistical criteria such as the Akaike Information Criterion and the Bayesian Information Criterion provide different methods for comparing models, while the \emph{Minimum Description Length}\index{Minimum Description Length} principle evaluates models through the information required to encode both the model and the data given the model. These approaches rely on different assumptions and do not necessarily select the same model. The fundamental problem is determining whether accuracy and complexity can be combined into a universal criterion or whether their relative importance depends on the objectives and assumptions of the investigation.

\emph{Problem of Explanatory Power}\index{Explanatory power}: Does a model that describes observations accurately and economically necessarily provide a satisfactory explanation? A model may achieve excellent predictive performance through complex computational procedures without revealing the mechanisms or relationships responsible for the observed phenomena. Conversely, a relatively simple model may provide substantial explanatory insight while making only approximate predictions. Proposed criteria for explanatory power include causal relevance, generality, unification, and the capacity to support counterfactual reasoning. These properties are not necessarily determined by empirical accuracy or descriptive complexity. The philosophical problem is establishing whether explanatory power can be evaluated independently of these quantities and whether a general measure of model adequacy can accommodate both predictive and explanatory objectives.

The evaluation of scientific models therefore involves several potentially competing criteria. Empirical accuracy, predictive reliability, simplicity, descriptive complexity, and explanatory power capture different aspects of model quality, and improvements in one dimension need not imply improvements in the others. The central philosophical challenge is to determine which properties are essential to scientific adequacy, how they can be measured or compared, and under what conditions trade-offs between them can be justified.

%
% Scientific Justification and Progress
%

\section{Scientific Justification and Progress}
\label{sec:scientific_justification}

\emph{Scientific justification}\index{Scientific justification} concerns the grounds on which scientific hypotheses are accepted, compared, revised, or rejected. Unlike mathematical theorems established through deductive proof, empirical scientific hypotheses generally depend on evidence that provides provisional and potentially revisable support. Moreover, different hypotheses may account for the same observations, and the criteria used to evaluate them may lead to conflicting judgments. These difficulties raise fundamental questions about the rational justification of scientific knowledge and the conditions under which scientific inquiry can be considered progressive.

\emph{Problem of Induction}\index{Problem of induction}: How can general scientific claims be justified by a finite collection of observations? \emph{Deductive reasoning}\index{Deduction} establishes conclusions that necessarily follow from their premises, whereas \emph{inductive reasoning}\index{Induction} extends beyond the information logically entailed by the available evidence. Hume's problem of induction concerns the justification of such inferences, particularly the assumption that regularities observed in the past will continue to hold in unobserved cases. This assumption cannot be established by deduction alone, while attempts to justify it inductively risk circular reasoning. Proposed responses include probabilistic accounts of induction, pragmatic justifications, and approaches that reject induction as a necessary foundation of science. The fundamental problem is determining whether general scientific claims can be rationally supported by limited evidence and, if so, under what conditions.

\emph{Problem of Confirmation and Falsification}\index{Confirmation}\index{Falsification}: What constitutes evidence for or against a scientific hypothesis? Confirmation theories investigate the conditions under which observations increase support for a hypothesis, whereas Popper's falsificationism emphasizes the importance of exposing scientific conjectures to potentially refuting tests. The \emph{hypothetico-deductive method}\index{Hypothetico-deductive method} evaluates hypotheses by deriving observational consequences and comparing them with empirical evidence. However, agreement with predictions does not establish the truth of a hypothesis, since alternative hypotheses may generate the same predictions. Conversely, an unsuccessful prediction need not uniquely falsify the hypothesis being tested, because the prediction may depend on auxiliary assumptions or experimental conditions. The problem is establishing rational criteria for evaluating positive and negative evidence without assuming that empirical tests provide conclusive justification.

\emph{Problem of Underdetermination}\index{Underdetermination}: Can scientific evidence uniquely determine which hypothesis should be accepted? The \emph{underdetermination of theory by evidence} arises when distinct, potentially incompatible hypotheses are consistent with the same observations. In particular, the \emph{Duhem--Quine thesis}\index{Duhem--Quine thesis} emphasizes that scientific hypotheses are generally tested together with auxiliary assumptions, making it difficult to determine which component should be revised when a prediction fails. Furthermore, different theories may be observationally equivalent within the available domain of investigation. This raises the question of whether empirical evidence alone can establish a preference between competing hypotheses or whether additional criteria, such as simplicity, explanatory power, or coherence, are required.

\emph{Problem of Probabilistic Justification}\index{Probabilistic justification}: Can the degree of support for a scientific hypothesis be quantified? Probabilistic approaches attempt to represent uncertainty and the evidential support of hypotheses through numerical quantities. In particular, \emph{Bayesian confirmation theory}\index{Bayesian confirmation theory} uses Bayes' theorem to update the probability assigned to a hypothesis in light of new evidence. This provides a formal method for combining prior assumptions with observational information. However, the resulting probabilities depend on the prior distribution, the likelihood model, and the hypotheses under consideration. Different assumptions may therefore produce different evaluations of the same evidence, while hypotheses not represented in the model cannot be directly assessed until the model is extended. The philosophical problem is determining whether probabilistic support provides an adequate account of scientific justification and whether the resulting degrees of belief can be interpreted objectively.

\emph{Problem of Theory Choice}\index{Theory choice}: How should scientists select among competing hypotheses that account for the available evidence? Proposed criteria include empirical adequacy, predictive accuracy, simplicity, explanatory power, generality, and coherence with established knowledge. The approach known as \emph{inference to the best explanation}\index{Inference to the best explanation} favors the hypothesis that provides the most satisfactory explanation of the evidence. However, these criteria may conflict: one hypothesis may provide better predictions, while another is simpler or offers greater explanatory scope. Moreover, the best available hypothesis need not be true, since superior alternatives may not yet have been formulated. The central problem is determining whether theory choice can be governed by universal criteria, how competing criteria should be balanced, and whether a uniquely preferred hypothesis can always be identified.

\emph{Problem of Scientific Acceptance and Revision}\index{Scientific acceptance}\index{Belief revision}: Under what conditions should a scientific hypothesis be accepted, modified, or rejected? Scientific inquiry involves revising existing commitments as new evidence becomes available, but such revisions need not consist simply of replacing a false hypothesis with a true one. A hypothesis may be retained despite anomalies, modified to accommodate additional observations, or replaced by an alternative with different assumptions and explanatory scope. Formal theories of \emph{belief revision}\index{Belief revision} investigate principles for incorporating new information while preserving consistency and minimizing unnecessary changes to existing beliefs. However, the rationality of revision also depends on the reliability of evidence, the structure of background assumptions, and the availability of alternatives. The problem is establishing criteria for rational scientific change without requiring that every new observation lead to the immediate rejection of existing knowledge.

\emph{Problem of Incommensurability}\index{Incommensurability}: Can scientific theories developed within different conceptual frameworks be evaluated according to common standards? Kuhn's account of scientific revolutions emphasizes that competing paradigms may differ in their concepts, explanatory objectives, and standards of evaluation. Such differences can make direct comparisons difficult because terms, observations, and criteria of success may not retain identical meanings across theoretical frameworks. However, incommensurability does not necessarily imply complete incomparability or the absence of rational theory choice. The philosophical problem is determining which properties must remain comparable when scientific frameworks change, whether common standards of evaluation can be established, and how rational comparisons can be made when the objects or criteria of investigation are themselves revised.

\emph{Problem of Scientific Progress}\index{Scientific progress}: What constitutes progress in scientific knowledge? A traditional cumulative conception regards progress as the accumulation of established facts and increasingly accurate theories. However, scientific revolutions may involve abandoning previously accepted concepts, changing explanatory frameworks, or redefining the phenomena under investigation. Alternative philosophical accounts characterize progress through increasing empirical success, greater explanatory or predictive power, improved problem-solving capacity, or closer approximation to truth. For example, Laudan associates progress with improvements in the ability of scientific theories and research traditions to solve empirical and conceptual problems, while truthlikeness approaches associate progress with movement toward a more accurate account of reality. The \emph{epistemic account of scientific progress}\index{Epistemic account of scientific progress}, defended by Bird, identifies progress with the accumulation of scientific knowledge. In contrast, the \emph{noetic account}\index{Noetic account of scientific progress}, developed by Dellsén, emphasizes improvements in scientific understanding, particularly the ability to correctly explain or predict aspects of the world. These approaches may yield different assessments of scientific developments, since acquiring additional knowledge does not necessarily produce greater understanding, and improvements in understanding need not correspond to the accumulation of new established facts. These perspectives need not agree about which scientific developments constitute improvements. The central problem is establishing the conditions under which changes in scientific knowledge can be regarded as progress rather than merely as changes in accepted theories.

\emph{Problem of Quantifying Scientific Progress}\index{Quantification of scientific progress}: Can scientific progress be measured, and can different states of scientific knowledge be objectively compared? A quantitative account requires identifying measurable properties of scientific knowledge and determining how changes in those properties correspond to improvement. Possible candidates include predictive accuracy, explanatory scope, problem-solving effectiveness, information content, and degrees of truthlikeness. However, these quantities capture different aspects of scientific achievement and may not vary together. A new hypothesis may improve empirical accuracy while increasing complexity, or expand explanatory scope while sacrificing precision in particular domains. Furthermore, comparisons may become difficult when scientific objectives, conceptual frameworks, or the entities under investigation change. The problem is determining whether progress admits a general quantitative characterization, whether different dimensions of improvement can be combined into a single measure, and whether scientific progress must always be understood as a cumulative or monotonic process.

These problems reveal that scientific justification and scientific progress are closely related but distinct. Justification concerns the rational grounds for accepting or preferring hypotheses, whereas progress concerns whether changes in scientific knowledge constitute genuine improvements. A hypothesis may be preferable to its competitors under particular criteria without establishing that knowledge has advanced in every relevant respect. The fundamental philosophical challenge is therefore to characterize the conditions under which scientific inquiry produces better knowledge and determine whether such improvements can be evaluated through objective, general, and potentially quantitative criteria.

%
% Scientific Discovery and the Limits of Knowledge
%

\section{Scientific Discovery and the Limits of Knowledge}
\label{sec:scientific_discovery}

\emph{Scientific discovery}\index{Scientific discovery} concerns the processes through which new knowledge, hypotheses, concepts, and entities become subjects of scientific inquiry. Discoveries may arise from systematic investigation, unexpected observations, theoretical reasoning, or the introduction of entirely new conceptual frameworks. However, discovering new knowledge is not equivalent to justifying it, and recognizing the existence of an unanswered question differs from identifying a question that had not previously been considered. Furthermore, scientific inquiry may encounter limitations arising from available methods, computational capabilities, or fundamental epistemological constraints. These considerations raise several philosophical problems concerning the possibility of discovering the unknown and determining the ultimate limits of scientific knowledge.

\emph{Problem of Scientific Discovery}\index{Problem of scientific discovery}: Can the generation of new scientific knowledge be characterized through general principles or systematic procedures? The traditional distinction between the \emph{context of discovery}\index{Context of discovery} and the \emph{context of justification}\index{Context of justification} separates the processes through which hypotheses are conceived from those through which they are evaluated. Some philosophical approaches emphasize the creative and historically contingent character of discovery, while others investigate whether it follows recognizable patterns of reasoning. Computational and methodological approaches further suggest that at least some discovery processes can be formalized. However, the existence of systematic discovery procedures for particular problems does not establish that a universal method of scientific discovery is possible. The fundamental question is whether discovery admits general principles, which aspects of it can be formalized, and whether such principles can extend beyond already identified scientific problems.

\emph{Problem of Abductive Discovery}\index{Abductive reasoning}: How can new scientific hypotheses be generated from existing observations? \emph{Abduction}\index{Abduction} is commonly understood as reasoning from observations to hypotheses that could explain them. When unexpected phenomena are encountered, researchers may formulate candidate explanations under which those phenomena would be less surprising. A distinction can be made between \emph{selective abduction}, which chooses among existing hypotheses, and \emph{creative abduction}, which introduces previously unconsidered hypotheses. However, abductive reasoning does not establish that the proposed explanation is correct, and several incompatible hypotheses may account for the same observations. Furthermore, identifying a phenomenon as surprising does not necessarily provide a procedure for generating an explanation. The problem is determining whether abductive reasoning can be formalized as a method of hypothesis generation and how genuinely novel hypotheses can arise from existing knowledge.

\emph{Problem of Heuristic Discovery}\index{Heuristic discovery}: Can scientific discovery be formulated as a systematic search through a space of possible hypotheses or solutions? Computational approaches to discovery often characterize scientific problems through a collection of possible states, operations that transform one state into another, and criteria for evaluating candidate solutions. \emph{Heuristic methods}\index{Heuristic methods} guide the search by prioritizing promising alternatives without guaranteeing that an optimal or satisfactory solution will be found. Such approaches can reproduce certain forms of scientific reasoning and assist in generating hypotheses. However, scientific problem spaces are frequently incomplete or poorly defined, and the relevant search operators, constraints, or evaluation criteria may themselves be unknown. A fundamental difficulty is that searching a predefined space cannot, without additional mechanisms, generate possibilities excluded by the space's original definition. The problem is determining whether systematic discovery requires a predefined space of possibilities and how such spaces can be expanded or transformed.

\emph{Problem of Conceptual Novelty}\index{Conceptual novelty}: How can scientific inquiry generate concepts, representations, or theories that cannot be expressed within an existing conceptual framework? Scientific discovery sometimes involves more than identifying new facts or combining familiar ideas; it may require changing the concepts through which phenomena are understood. Approaches based on \emph{analogical reasoning}\index{Analogical reasoning} and \emph{model-based reasoning}\index{Model-based reasoning} investigate how relationships from familiar domains can be transferred to unfamiliar situations or how models can be manipulated to generate new possibilities. Boden distinguishes between combinational, exploratory, and transformational creativity, the latter involving changes to the rules defining a conceptual space. Kuhn's account of scientific revolutions similarly emphasizes the possibility of substantial changes in scientific concepts and frameworks. The philosophical problem is determining how genuinely new concepts can arise from existing knowledge and whether conceptual transformations can be characterized through general operations or rules.

\emph{Problem of Discovering Unknown Entities}\index{Discovery of unknown entities}: How can scientific inquiry identify entities or phenomena whose existence has not previously been suspected? Some discoveries concern new properties of recognized entities, while others introduce previously unrecognized objects of investigation. The latter may arise through unexpected observations, the reinterpretation of existing evidence, or theoretical predictions suggesting the existence of new entities. Discoverability must nevertheless be distinguished from knowability: an entity may be knowable in principle without there being any available procedure for identifying it as a subject of investigation. Conversely, identifying a candidate entity does not establish that it exists or that its properties can be known. However, a novel observation does not necessarily indicate a novel entity, and a newly proposed entity may subsequently prove to be fictitious or a misinterpretation of existing phenomena. Furthermore, the criteria used to recognize potential discoveries depend on existing conceptual frameworks, which may prevent researchers from identifying phenomena that fall outside their expectations. The problem is determining how previously unrecognized entities can become identifiable, how genuine novelty can be distinguished from observational or conceptual error, and whether systematic methods for discovering unknown entities are possible.

\emph{Problem of Epistemic Ignorance}\index{Epistemic ignorance}: What distinguishes what is known from what is unknown, and can ignorance itself be recognized? A distinction is commonly made between \emph{known unknowns}\index{Known unknowns}, concerning questions or subjects whose unresolved status is recognized, and \emph{unknown unknowns}\index{Unknown unknowns}, concerning questions, entities, or possibilities that have not yet entered awareness. These categories are relative to a particular individual, community, domain of inquiry, and state of knowledge. Moreover, recognizing that a question is unanswered does not imply knowing its answer, while failing to recognize a question does not imply that it is unanswerable. A central difficulty is that unknown unknowns cannot generally be identified individually without thereby becoming recognized subjects of inquiry. The philosophical problem is determining whether the unknown can be characterized without presupposing complete knowledge of its contents, and whether a meaningful boundary can be established between recognized and unrecognized ignorance.

\emph{Problem of Knowability}\index{Problem of knowability}: Are all entities and truths knowable in principle, or are some fundamentally beyond the reach of knowledge? It is necessary to distinguish between limitations arising from the current state of scientific knowledge and limitations that would persist even under ideal conditions. An entity may be inaccessible using existing instruments or conceptual frameworks while remaining knowable through methods not yet developed. Conversely, some limits may follow from logical, mathematical, or physical constraints. Knowability may also depend on the epistemic capacities and admissible methods of a particular investigator or class of investigators. Philosophical discussions of knowability examine whether every truth could in principle become known and what assumptions such a claim requires. The problem is determining whether general criteria for knowability can be established, how relative limitations can be distinguished from absolute ones, and whether unknowability can itself be demonstrated.

\emph{Problem of Computational and Logical Limits}\index{Computational limits}\index{Logical limits of knowledge}: What limitations do logic and computation impose on the acquisition of scientific knowledge? Results from mathematical logic and computability theory establish fundamental restrictions on particular classes of formal procedures. For example, the undecidability of the halting problem demonstrates that no algorithm can correctly decide termination for every possible program and input. Similarly, Gödel's incompleteness theorems establish limitations on sufficiently expressive, consistent, effectively axiomatized formal systems. However, these results do not imply that every difficult scientific question is undecidable or that every truth inaccessible to a particular formal procedure is unknowable by all possible means. Their implications depend on the formal assumptions and the methods of inquiry under consideration. The philosophical problem is determining how mathematical limitations on formal systems and computational procedures constrain scientific knowledge, and whether these restrictions establish absolute epistemic boundaries or limitations relative to particular methods.

\emph{Problem of Perfect Knowledge}\index{Perfect knowledge}: Is complete knowledge of an entity or phenomenon possible, and under what conditions could it be recognized? An ideal of perfect knowledge may require a complete and accurate account of all relevant properties, relationships, and behavior of an entity. However, the meaning of completeness depends on the questions being considered, the level of detail required, and the standards adopted for scientific understanding. An account that is complete relative to one domain of inquiry may remain incomplete relative to another. Furthermore, even if complete knowledge exists, establishing that it has been attained may require ruling out unknown properties, alternative explanations, or future observations that could challenge the current account. This creates an important distinction between the existence of perfect knowledge, the possibility of acquiring it, and the ability to recognize its attainment. The philosophical problem is determining whether perfect knowledge admits a meaningful general characterization and whether scientific inquiry can establish that no further improvement is possible.

These problems illustrate two complementary aspects of scientific inquiry: the expansion of knowledge through the discovery of new entities, hypotheses, and concepts, and the recognition of limitations that may constrain such expansion. Scientific discovery need not proceed through a universal method, and the boundaries of knowledge may depend on both the available means of investigation and fundamental logical or physical constraints. The central philosophical challenge is to determine how the unknown can be meaningfully characterized, whether previously unrecognized problems can be systematically discovered, and under what conditions scientific inquiry can approach or attain complete understanding.

%
% Section: References
%

\section*{References}

\cite{chalmers2013thing}: Chalmers provides a clear and widely used introduction to the philosophy of science, discussing induction, falsifiability, paradigms, and the limits of scientific reasoning. It is ideal as a general reference for framing the main questions addressed in this chapter.

\cite{popper2014conjectures}: Popper introduces falsifiability as the demarcation criterion for science and develops his vision of knowledge growth through bold conjectures and refutations. This is central for your treatment of scientific progress and the demarcation problem.

\cite{van1980scientific}: Van Fraassen defends constructive empiricism, arguing that science aims for empirical adequacy rather than truth. This perspective supports your discussion on the limits of science and the interpretation of scientific representations.

\cite{suppes2002representation}: Suppes develops a structuralist account of scientific theories, emphasizing how scientific knowledge can be formalized in terms of models and invariants. This is directly relevant to your sections on scientific representation.

\cite{mosterin2016conceptos}: Mosterín explores the logical and conceptual structure of scientific theories, focusing on how concepts evolve and how theories are built. This supports your account of entities, theories, and the evolution of knowledge.

\cite{latour2013laboratory}: Latour's ethnographic study highlights the social processes behind the making of scientific knowledge. Including this reference enriches your treatment of science as a human practice and its limits.

\cite{nagel2014knowledge}: Nagel gives a concise overview of epistemological issues such as justification, belief, and truth. This is useful background for connecting philosophy of science with broader epistemology.

\ref{pirsig1999zen} contains an interesting review of the concept of science, the scientific method, and the role that technology plays in our society. The author proposes that the goal of science should be quality, although the concept of quality is left undefined, and how to reconcile the rational and romantic points of view in science. The book also contains some advice about which is the right state of mind to pursue a scientific problem, and how to deal with the inevitable failures.

% New References

% Edmund Gettier — Is Justified True Belief Knowledge? (1963).
\cite{gettier1963justified}: Gettier challenges the traditional definition of knowledge as justified true belief, demonstrating that truth and justification are not sufficient to establish knowledge.

% Luciano Floridi — The Philosophy of Information (2011).

\cite{floridi2011philosophy}: Floridi examines the philosophical foundations of information and its relationship to meaning, truth, and knowledge. His analysis is particularly relevant to the distinction between information-theoretic quantities and epistemic properties.

% Representing and Intervening — Ian Hacking (1983).

\cite{hacking1983representing}: Hacking examines scientific realism and the ontological status of theoretical entities, emphasizing the role of observation and experimental intervention in evaluating claims about their existence.

% Natural Categories and Human Kinds — Muhammad Ali Khalidi (2013).

\cite{khalidi2013natural}: Khalidi investigates the philosophical foundations of scientific classification and natural kinds, addressing whether scientific categories correspond to objective divisions in reality or reflect particular explanatory purposes.

% Modelling Nature: An Opinionated Introduction to Scientific Representation — Roman Frigg and James Nguyen (2020)

\cite{frigg2020modelling}: Frigg and Nguyen provide a systematic examination of the fundamental problems of scientific representation, the conditions that an adequate philosophical account should satisfy, and the principal competing approaches to representation.

% An Inferential Conception of Scientific Representation — Mauricio Suárez (2004)

\cite{suarez2004inferential}: Suárez develops an inferential conception of scientific representation based on directionality and surrogative reasoning, critically examining accounts based on similarity and structural isomorphism.

% Aspects of Scientific Explanation and Other Essays in the Philosophy of Science — Carl G. Hempel (1965)

\cite{hempel1965aspects}: Hempel develops the deductive-nomological and statistical approaches to scientific explanation, examining the role of general laws, logical inference, and the relationship between explanation and prediction.

% Making Things Happen: A Theory of Causal Explanation — James Woodward (2003) Oxford University Press.

\cite{woodward2003making}: Woodward develops an interventionist account of causal explanation, addressing causal relevance, counterfactual dependence, and the conditions under which scientific models support explanations of phenomena.

% How to Tell When Simpler, More Unified, or Less Ad Hoc Theories Will Provide More Accurate Predictions — Malcolm Forster and Elliott Sober (1994)

\cite{forster1994simpler}: Forster and Sober examine the epistemic justification of simplicity, unification, and parsimony, investigating the conditions under which simpler models may provide more accurate predictions.

% The Minimum Description Length Principle — Peter D. Grünwald (2007)

\cite{grunwald2007minimum}: Grunwald develops the Minimum Description Length principle as a framework for statistical inference and model selection, addressing the relationship between descriptive complexity, empirical fit, and overfitting.

% Progress and Its Problems: Towards a Theory of Scientific Growth — Larry Laudan (1977)

\cite{laudan1977progress}: Laudan examines the philosophical foundations of scientific progress, proposing problem-solving effectiveness as a criterion for evaluating scientific theories and research traditions.

% Is Science Progressive? — Ilkka Niiniluoto (1984)

\cite{niiniluoto1984science}: Niiniluoto investigates the relationship between scientific progress, truth, and truthlikeness, examining whether improvements in scientific knowledge can be objectively characterized and compared.

% Additonal references

\cite{kukla2023knowing}: Kukla examines objectual knowledge as an epistemic relationship distinct from propositional knowledge, investigating what it means to know an object or place.

\cite{piedrahita2025capacity}: Piedrahita develops a capacity-based conception of ignorance, challenging accounts that define ignorance exclusively through the absence of knowledge or true belief.

\cite{simion2026ignorance}: Simion, Kelp, and Mace examine the epistemology and semantics of ignorance, including propositional and objectual ignorance and the possibility of characterizing degrees of ignorance.

\cite{bird2007progress}: Bird develops an epistemic account of scientific progress, according to which scientific progress consists in the accumulation of scientific knowledge.

\cite{dellsen2016progress}: Dellsén develops a noetic account of scientific progress based on improvements in scientific understanding, challenging the identification of progress exclusively with the accumulation of knowledge.
